%\documentclass[a4pape, 14pt]{extarticle} % расширенный класс статьи с возможностью указать 14 шрифт
%\documentclass[a4pape, 11pt]{article}
\documentclass[oneside, final, 14pt]{extarticle}
\usepackage{graphicx}  % поддержка .eps-графики
\usepackage[utf8]{inputenc} % кодировка в которой набран текст 
\usepackage[T2C]{fontenc} % поддержка кирилицы % ещё можно T1 T2B T2A T2
\usepackage[english,russian]{babel} % переносы и типографские правила для русского
\usepackage{indentfirst}  % красная строка
\usepackage{listings}  % оформление листингов программ
\usepackage{amssymb,amsfonts,amsmath,mathtext} % ядро для научной статьи
\usepackage{physics}
\usepackage{mathtools}
\usepackage{braket}
\usepackage{float} % для рисунков
\usepackage{array} % для m{2cm} в таблицах
\usepackage{blindtext} % случайный текст для заглушки

\usepackage{titlesec}
%\linespread{1.3}
%\usepackage{setspace}\onehalfspacing
%\usepackage{vmargin} отступы из книги столярова
%\setpapersize{A4}
%\setmarginsrb{2.5cm}{2cm}{1.5cm}{2cm}{0pt}{0mm}{0pt}{13mm}


%\usepackage{sectsty} % не работает с extarticle. Должно было уменьшить размер заголовков
%\sectionfont{\fontsize{14}{14}\selectfont}

\usepackage{titlesec} % https://ru.overleaf.com/learn/latex/Sections_and_chapters % пакет для изменения заголовков разделов
\titleformat{\section}
{\normalfont\fontsize{14}{14}\bfseries}{\thesection}{1em}{}

\usepackage{pdfpages} % для вставки пдф файлов

\usepackage{pgfplots} % построение графиков
% We will externalize the figures
%\usepgfplotslibrary{external}
%\tikzexternalize
%\pgfplotsset{width=15cm,compat=1.5}

%\sloppy
%\fussy

\usepackage[left=2cm, right=2cm, top=1.5cm, bottom=1.5cm]{geometry} % установка полей

\newcommand{\hup}{Принцип неопределенности Гейзенберга}
\newcommand{\qe}{Квантовая запутанность}

%\usepackage{citehack}  % https://www.opennet.ru/docs/RUS/cyr_howto/ch08s02.html
%\usepackage[bibencoding=auto,backend=biber,babel=other]{biblatex}
%\usepackage{csquotes} 


% https://ru.overleaf.com/learn/latex/Pgfplots_package  - для построения графиков
% https://ru.overleaf.com/learn/latex/Tables - построение таблиц
% https://ru.overleaf.com/learn/latex/Commands - обьявление команд
\graphicspath{{./images/}} %путь к рисункам
\frenchspacing % длина пробелов после пунктуации
\pagestyle{plain}
\selectlanguage{russian}
\bibliographystyle{gost780} % gost780 - сортировка по порядку следования, gost780s - по алфавиту кажется. gost780u не завелся

\begin{document}

\begin{titlepage}
	\centerline{\includegraphics[width=2cm]{dstu-logo}}
	\vfill
	%\Large
	\centerline{МИНЕСТЕРСТВО НАУКИ И ВЫСШЕГО ОБРАЗОВАНИЯ}
	\centerline{РОССИЙСКОЙ ФЕДЕРАЦИИ}
	\vfill
	\centerline{\bf ФЕДЕРАЛЬНОЕ ГОСУДАРСТВЕННОЕ БЮДЖЕТНОЕ}
	\centerline{\bf ОБРАЗОВАТЕЛЬНОЕ УЧРЕЖДЕНИЕ ВЫСШЕГО ОБРАЗОВАНИЯ}
	\centerline{\bf «ДОНСКОЙ ГОСУДАРСТВЕННЫЙ ТЕХНИЧЕСКИЙ УНИВЕРСИТЕТ»}
	\centerline{\bf (ДГТУ)}
	\normalsize
	\vfill\vfill
	\centerline{Факультет «Отдел аспирантуры и докторантуры»}
	\centerline{Кафедра «Кибербезопасность информационных систем»}
	\vfill
	%\rightline{Место для подписи завкафедрой}
	\vfill
	%\Large{\centerline{\bf ПОРТФОЛИО АСПИРАНТА}}
	\centerline{План-проспект реферата по дисциплине «История и философия науки»}
%	\centerline{«История и философия науки»}
	\centerline{на тему «История и философия криптографии»}
	\vfill
	\vfill
	\vfill
	\vfill
	\vfill
	\vfill
	\vfill
	\rightline{Выполнил:}
	\rightline{аспирант 1го курса}
	\rightline{по направлению подготовки}
	\rightline{09.06.01 «Информатика и вычислительная техника»}% номер зачетной книжки: 2142748
	\rightline{Галов К.А.}
	\rightline{Проверил: д.ф.н., доцент Тазаян А.Б.} % Тазаян Араван Бабкенович
	\vfill
	\vfill
	
	\centerline{г. Ростов-на-Дону}
	\centerline{2022 г.}
	
\end{titlepage}
\setcounter{page}{2}

\section*{Введение}
Определение криптографии. Актуальность этого раздела науки. Цель криптографии.

\section{Идеалы и нормы научности в криптографии}
Принципы написания научных работ по криптографии. Идеализация участников канала связи.
% обычные принципы + то, что злоумышленник ограничен только законами физики.

Литература: \cite{phil2016nesm, scarani2004quantum}

\section{Древнейшие шифры}
Шифрование в древнем Риме. Особенности первых методов шифрования. 

Литература: \cite{senthil2013modern}

\section{Средневековые шифры}
Шифры времен средневековья. Шифр Виженера. 

Литература: \cite{omran2011cryptanalytic}

\section{Криптография времен второй мировой войны}
Шифровальная машина Энигма. Работы Алана Тьюринга. Шифровальная техника СССР.

Литература: \cite{d2006sov, g2005n, hodges2014alan}


\section{Проблемы современной криптографии}
Основы современной криптографии. Теоретическая ненадежность современной криптографии.

Литература: \cite{bellare2005introduction}

\section{Квантовая криптография}
Основы квантовой криптографии. Надежность квантовой криптографии.

Литература: \cite{scarani2004quantum, ch1984quantum, lo2005decoy}

\section{Заключение}
Общий вывод об истории криптографии, как науки и обзор дальнейших перспектив развития.

\bibliography{refs}

\clearpage
\pagenumbering{gobble}
\section*{Отзыв научного руководителя на используемую литературу}
Я, Черкесова Лариса Владимировна, научный руководитель аспиранта первого года обучения Галова Кирилла Алексеевича подтверждаю, что данная тема соответствует специальности 09.06.01 «Информатика и вычислительная техника» и направлению исследований аспиранта.
\bigskip

\rightline{проф. каф. «КБИС», д.ф.-м.н. \hspace{0.5cm} \underline{\hspace{3cm}} \hspace{0.5cm} Черкесова Л. В}

\end{document}